% Image credits for the photographs and AI illustrations of Book 3 (University, Year 1).
% Machine-readable license records live in images/book3/CREDITS.md; keep the
% two files in sync. The Book 3 writer owns this file: add one \item per
% photograph (inside the itemize, by chapter) as photographs land.
\chapter*{Image Credits}

The drawings, maps and charts in this book are original; maps and charts
drawn from published data name their source in the caption or in this
list. Photographs are used under their respective free licenses, with
thanks to their authors; full source links and license URLs for every
photograph are listed in \texttt{images/book3/CREDITS.md} in the book's
repository.

\bigskip
Alongside the photographs and the original drawings, some figures are
AI-generated illustrations, produced with OpenAI image generation from
prompts written by the book's editors, and reviewed for scientific
accuracy before inclusion. They show generic scenes only, never a real
place, a real person or a real specimen. The full prompt for every
generated image is recorded in \texttt{images/book3/ai/PROMPTS.md} in the
repository.

\bigskip
\noindent Photographs and data, by chapter:
\begin{itemize}\small
  \item Ch.~1 (spherical astronomy): star trails above the Paranal Observatory, R.~Wesson/ESO, CC BY 4.0.
  \item Ch.~2 (telescopes, detectors and photometry): six segments of the primary mirror of the James Webb Space Telescope before cryogenic testing, NASA/MSFC/David Higginbotham, public domain.
  \item Ch.~3 (comparative planetology): the terrestrial planets to scale, NASA; Venus from Magellan radar, NASA/JPL; Mars from the Viking 1 orbiter, NASA/USGS; Europa, NASA/JPL/DLR; Io, NASA/JPL/University of Arizona; all public domain.
  \item Ch.~4 (geological maps): William Smith, \emph{A Delineation of the Strata of England and Wales, with part of Scotland}, 1815; public domain (Wikimedia Commons).
  \item Ch.~5 (cross-sections): Hutton's unconformity at Siccar Point, Pawe\l{}S, CC BY-SA 4.0 (Wikimedia Commons).
  \item Ch.~6 (field and laboratory methods): a biotite granite in thin section between crossed polars, Darla Sondrol, CC0 (Wikimedia Commons).
  \item Ch.~7 (rock-forming minerals): olivine, quartz, biotite, almandine and hornblende specimens, Robert M.~Lavinsky, CC BY-SA 3.0; labradorite, Twyla Baker, CC0; calcite cleavage rhombohedra, Giovanna Canu and Eva Santini, CC BY 4.0 (all Wikimedia Commons).
  \item Ch.~8 (optical mineralogy): a gabbro in thin section with one polar and between crossed polars, Chmee2, CC BY-SA 3.0; Auguste Michel-L\'evy as a student, public domain (Wikimedia Commons). The interference colours are computed from the CIE 1931 colour-matching functions and CIE standard illuminant D65, \copyright{} CIE, CC BY-SA 4.0.
  \item Ch.~9 (the layered Earth): Andrija Mohorovi\v{c}i\'c, public domain; Inge Lehmann in 1932, photograph by Even Neuhaus, public domain (Wikimedia Commons). PREM: A.~M.~Dziewonski and D.~L.~Anderson (1981), as tabulated in ObsPy.
  \item Ch.~10 (the Earth's heat budget): global heat flow from J.~H.~Davies and D.~R.~Davies (2010), \emph{Solid Earth}, CC BY 3.0; mantle adiabat from T.~Katsura (2022), \emph{JGR Solid Earth}, CC BY.
  \item Ch.~11 (isostasy and gravity anomalies): the Vening Meinesz pendulum apparatus, Broodt, CC BY-SA 4.0 (Wikimedia Commons). Elevations from ETOPO1 (NOAA National Centers for Environmental Information), sampled at one degree.
  \item Ch.~12 (plate kinematics on a sphere): a permanent satellite-positioning station, Alpsdake, CC BY-SA 4.0 (Wikimedia Commons). Plate boundaries from PB2002 (P.~Bird, 2003; GeoJSON by H.~Ahlenius, ODC-By 1.0); coastlines from Natural Earth (public domain).
  \item Ch.~13 (rifts, ridges and passive margins): open fissures at \TH{}ingvellir, Iceland, Olga Ernst, CC BY-SA 4.0; pillow lavas on the Gal\'apagos Rift, NOAA Okeanos Explorer Program, CC BY 2.0 (both Wikimedia Commons). Depths from ETOPO1 (NOAA National Centers for Environmental Information) and ocean-floor ages from M.~Seton and others (2020), EarthByte, CC BY 4.0.
  \item Ch.~14 (active margins): Sarychev Peak erupting, seen from the International Space Station, NASA, public domain (Wikimedia Commons). Earthquakes from the USGS ANSS Comprehensive Catalog (public domain); volcanoes from the Smithsonian Institution Global Volcanism Program; plate boundaries from PB2002 (P.~Bird, 2003; GeoJSON by H.~Ahlenius, ODC-By 1.0); coasts from Natural Earth (public domain); ocean-floor ages from M.~Seton and others (2020), EarthByte, CC BY 4.0.
  \item Ch.~15 (hotspots and intraplate volcanism): a lava fountain at Pu\textquoteleft u \textquoteleft O\textquoteleft o, Kilauea, G.~E.~Ulrich, U.S. Geological Survey, public domain; the Western Ghats at Matheran, Nicholas (Nichalp), CC BY-SA 2.5 (both Wikimedia Commons). Hotspot positions, track azimuths and rates from Table~1 of the electronic supplement of W.~J.~Morgan and J.~P.~Morgan (2007); depths from ETOPO1 (NOAA National Centers for Environmental Information); plate boundaries from PB2002 (P.~Bird, 2003; GeoJSON by H.~Ahlenius, ODC-By 1.0); coasts from Natural Earth (public domain).
  \item Ch.~16 (stress, strain and rheology): a high-pressure triaxial cell for measuring rock strength, Darryl Peroni, CSIRO, CC BY 3.0 (Wikimedia Commons). Dislocation-creep parameters of wet quartzite (Gleason and Tullis, 1995) and dry olivine (Hirth and Kohlstedt, 2003) as compiled by M.~A.~Lopez-Sanchez in \emph{strength\_envelopes} (MPL-2.0).
  \item Ch.~17 (faults, folds and fabrics): a slickenside at Corona Heights, San Francisco, Dextersinister, CC BY-SA 4.0; fault breccia of the Keystone Thrust, Nevada, Jstuby, public domain; chevron folds at Millook, Cornwall, Smalljim, CC BY 3.0; an anticline with axial-planar cleavage near Potes, Cantabria, Mikenorton, CC BY-SA 3.0 (all Wikimedia Commons).
  \item Ch.~18 (the earthquake source): land and fences separated by the 1906 rupture in Marin County, \emph{Popular Science Monthly} vol.~69 (1906), public domain; Wallace Creek offset across the San Andreas fault, R.~E.~Wallace, U.S. Geological Survey, public domain (both Wikimedia Commons). Earthquakes from the USGS ANSS Comprehensive Catalog (public domain); the focal mechanism of the 2011 Tohoku earthquake from the USGS W-phase moment tensor (public domain).
  \item Ch.~19 (weathering and clays): rotten granite and granite sand at Targassonne, Pyrenees, Pinpin, CC~BY~2.5; bauxite on kaolinitic sandstone at Weipa, Australia, Werner Schellmann, CC~BY-SA~2.5; kaolinite crystals under the scanning electron microscope, U.S. Geological Survey, public domain (all Wikimedia Commons). Ionic radii from R.~D.~Shannon (1976).
  \item Ch.~20 (erosion and transport): the flood plume of the Mississippi from the MODIS instrument on Terra, 4~March 2018, NASA Earth Observatory image by Joshua Stevens, public domain (Wikimedia Commons).
  \item Ch.~21 (clastic sediments and their environments): cross-bedding in the Aztec Sandstone, Nevada, Chip Carroon, Bureau of Land Management, CC~BY-SA~2.0; the Lena River delta from Landsat~7, U.S. Geological Survey, public domain; a turbidite with the divisions of the Bouma sequence, Becke-Oese, Jo Weber, CC~BY-SA~3.0 (all Wikimedia Commons).
  \item Ch.~22 (carbonate, siliceous and evaporite sediments): Bora Bora from the International Space Station, NASA Johnson Space Center, public domain; a cell of \emph{Emiliania huxleyi}, Alison R.~Taylor, University of North Carolina Wilmington Microscopy Facility, CC~BY~2.5; \emph{Orbulina universa}, H.~Takagi, K.~Kimoto, T.~Fujiki, H.~Saito, C.~Schmidt, M.~Kucera and K.~Moriya, CC~BY~4.0; diatoms of McMurdo Sound, Gordon T.~Taylor, Stony Brook University, NOAA Corps Collection, public domain (all Wikimedia Commons).
  \item Ch.~23 (diagenesis and sedimentary basins): quartz overgrowths in the Potsdam Sandstone, M.~C.~Rygel, CC~BY-SA~4.0 (Wikimedia Commons). Sea-level stack of R.~M.~Spratt and L.~E.~Lisiecki (2016), NOAA NCEI World Data Service for Paleoclimatology; compaction parameters of J.~G.~Sclater and P.~A.~F.~Christie (1980) as tabulated in pyBacktrack (EarthByte).
  \item Ch.~24 (the stratigraphic scale): golden spike of the Selandian GSSP at Zumaia, Josugoni, CC~BY-SA~4.0 (Wikimedia Commons); the Cretaceous--Paleogene boundary at Trinidad Lake State Park, Nationalparks, CC~BY-SA~2.5 (Wikimedia Commons). Base ages and colours of the International Chronostratigraphic Chart, International Commission on Stratigraphy, CC~BY~4.0.
  \item Ch.~25 (radiometric dating methods): a zircon grain in cathodoluminescence, Emmanuel Roquette, CC~BY-SA~4.0; a thermal ionisation mass spectrometer, Radiogenic, CC0 (both Wikimedia Commons). The IntCal20 calibration curve of P.~J.~Reimer and others (2020), \emph{Radiocarbon} 62, CC~BY~4.0, as distributed with the rintcal package (M.~Blaauw).
  \item Ch.~26 (structure of the atmosphere and the ocean): retrieving a CTD and rosette in the Bering Strait, NASA Goddard Space Flight Center, public domain (Wikimedia Commons). Temperature and salinity of the World Ocean Atlas 2023 (NOAA National Centers for Environmental Information).
  \item Ch.~27 (radiation and heat transport): the sunlit Earth from the Deep Space Climate Observatory, NASA Goddard Space Flight Center, public domain (Wikimedia Commons). Radiation, surface fluxes and precipitation of the NCEP-NCAR Reanalysis~1 (E.~Kalnay and others, 1996), data provided by the NOAA PSL, Boulder, Colorado, USA, from their website at https://psl.noaa.gov.
  \item Ch.~28 (the general circulation of the atmosphere): an extratropical cyclone in the Gulf of Alaska, NASA MODIS/LANCE, public domain; the intertropical convergence zone from GOES-11 cloud data, NASA Earth Observatory, public domain (both Wikimedia Commons). Winds and sea-level pressure of the NCEP-NCAR Reanalysis~1 (E.~Kalnay and others, 1996), data provided by the NOAA PSL, Boulder, Colorado, USA, from their website at https://psl.noaa.gov.
  \item Ch.~29 (the circulation of the ocean): the Gulf Stream seen through sea-surface temperature, NOAA NESDIS, public domain; a deep profiling float being deployed, unidentified NOAA photographer, public domain (both Wikimedia Commons).
\end{itemize}
